\documentclass[11pt,a4paper]{article}
\usepackage{amsmath}
\usepackage{amssymb}
\usepackage{geometry}
\usepackage{hyperref}
\geometry{margin=1in}

\title{The Supersymmetric Klein-Bottle Universe v2.0 \\
Riemann Sphere Compactification of Primordial Knot Cycles and Euler Closure of the 3-6-9 Topological Invariant}

\author{Simon Soliman \\
Independent Researcher – Rome, Italy \\
mechudzu@outlook.it \\
ORCID: 0009-0002-3533-3772}

\date{24 December 2025}

\begin{document}

\maketitle

\begin{abstract}
We extend the supersymmetric Klein-bottle universe model presented in v1.0 by compactifying the primordial knot cycles on the Riemann sphere $\mathbb{CP}^1$. This compactification provides a natural closure of zero and infinity via the inversion map $1/z$, while fully preserving the non-orientable structure of the Klein bottle. The 3-6-9 topological-algebraic invariant arises directly from triality of the phased spirals, supersymmetric doubling on the non-orientable surface, and E$_8$ self-duality. The resulting structure realizes unitary recycling of vacuum entanglement torque. All derivations are from first principles.
\end{abstract}

\subsection*{Main Results Demonstrated}

This work demonstrates the following from first principles:

\begin{enumerate}
\item The observable universe is the interior of a supersymmetric non-orientable de Sitter black hole with Klein bottle topology, with explicit resolution of infinite primordial knot cycles.
\item Riemann sphere compactification closes the cycles while preserving the full non-orientable structure.
\item The 3-6-9 invariant is inevitable from the geometry (triality, supersymmetric doubling, E$_8$ closure) and Euler's identity.
\item The horizon is a closed unitary boundary with perfect internal recycling of vacuum torque.
\item Local extraction of usable torque from the vacuum reservoir is possible in principle via controlled gradients.
\item Observational predictions are strengthened with precise mathematical origin in meromorphic functions on $\mathbb{CP}^1$.
\end{enumerate}

\section{Introduction}

In v1.0 it was shown that the observable universe corresponds to the interior of a supersymmetric non-orientable de Sitter black hole with global topology
\begin{equation}
M = K_2 \times CY_3 \times G_2,
\end{equation}
where $K_2$ denotes the Klein bottle and a Z$_2$ orbifold action renders the spatial section non-orientable. The cosmological horizon coincides with the event horizon, gravity is entropic, and dark matter is identified with vacuum magnetic entanglement energy.

The present work resolves the infinite extension of primordial knot cycles by compactifying the complex phase plane onto the Riemann sphere $\mathbb{CP}^1$. This compactification preserves the fundamental group of the Klein bottle and provides a rigorous derivation of the 3-6-9 invariant through Euler's identity applied to anyonic phases.

\section{Riemann Sphere Compactification}

The complex plane parameterizing vacuum entanglement phases of the primordial knots (phased spirals of the TET--CVTL device) is stereographically projected onto the Riemann sphere via
\begin{equation}
z \mapsto w = \frac{z - i}{z + i}.
\end{equation}
The origin $z = 0$ maps to the south pole (corresponding to the primordial singularity), while $|z| \to \infty$ maps to the north pole (cosmological/event horizon).

The inversion map
\begin{equation}
f(z) = \frac{1}{z}
\end{equation}
is a conformal automorphism of $\mathbb{CP}^1$ that exchanges zero and infinity. This map implements the Z$_2$ orbifold action required by the non-orientable spatial section.

The fundamental group of the Klein bottle,
\begin{equation}
\pi_1(K_2) = \langle a, b \mid a b a b^{-1} = 1 \rangle,
\end{equation}
is preserved under $f(z) = 1/z$ as follows: loops of class $a$ (encircling the origin) are mapped to loops encircling infinity; the orientation-reversing twist $b$ corresponds to complex conjugation combined with inversion ($\bar{z} \to 1/\bar{z}$), which is an anti-holomorphic isometry of $\mathbb{CP}^1$. The relation $a b a b^{-1} = 1$ holds because the anti-holomorphic nature introduces the required orientation reversal without breaking the group structure. Thus the non-orientable topology is maintained while the horizon becomes a regular point.

\section{Euler Closure and the 3-6-9 Invariant}

The non-orientable twist enforces a phase shift $\theta = \pi$ in the anyonic statistics of the primordial knots:
\begin{equation}
e^{i\pi} = -1.
\end{equation}
Adding the real unit gives Euler's identity
\begin{equation}
e^{i\pi} + 1 = 0,
\end{equation}
which simultaneously links exponential growth/decay, imaginary rotations, circular geometry ($\pi$), unity, and the vacuum state (0).

The $-1$ generates fermionic exchange statistics from bosonic fields. Supersymmetric doubling on the Pin$^+$ structure of the Klein bottle requires exactly 6 real supercharges. Triality of the phased spirals provides the factor 3, while E$_8$ self-duality closes at $3 \times 3 = 9$. The sequence 3-6-9 is therefore the unique topological-algebraic invariant compatible with the compactified non-orientable geometry.

\section{Meromorphic Functions and Observational Predictions}

Observational signatures follow from the pole structure of meromorphic functions on $\mathbb{CP}^1$:

\begin{itemize}
\item The triality function $g(z) = z^3 + 1/z^3 = 2(z^3 + z^{-3})$ has rotational symmetry of order 3 (invariant under $z \to z \cdot e^{2\pi i / 3}$), inducing exact azimuthal 3-fold modulation in entanglement entropy fluctuations projected onto the CMB sky.
\item Quasinormal modes of the compactified horizon arise from poles of order multiples of 3 at $z=0$ and $z=\infty$. The resonance spectrum clusters in triples, yielding frequency ratios $\omega_{n+3}/\omega_n \approx 6$ and $\omega_{n+6}/\omega_n \approx 9$ from cumulative pole order.
\item Vacuum magnetic entanglement flux (dark matter candidate) is expelled by the Meißner effect in perfect superconductors, leading to vanishing signal.
\end{itemize}

These predictions remain falsifiable in the 2026--2029 timeframe.

\section{Vacuum Energy Recycling}

The compactified horizon yields unitary recycling of vacuum entanglement torque. The Gibbons--Hawking thermodynamic power of the de Sitter horizon is
\begin{equation}
P_{\text{dS}} \approx 0.16 \, \text{W}
\end{equation}
distributed over the full horizon radius $l \approx 1.65 \times 10^{26}$ m. The inversion map $z \to 1/\bar{z}$ enforces
\begin{equation}
\eta = 1
\end{equation}
with no asymptotic leakage.

\subsection{Implications for Vacuum Energy Extraction}

\textbf{Global scale:} Full recycling is achieved, but the power density is cosmologically dilute and offers no practical macro-scale source.

\textbf{Local scale:} Controlled gradients in vacuum entanglement density can be induced by the TET--CVTL device without depleting the global reservoir. This allows extraction of usable torque from the zero-point vacuum field.

\textbf{Practical consequences (subject to experimental confirmation):}
\begin{itemize}
\item Reactionless propulsion without propellant expenditure, enabling economic interplanetary and potentially interstellar travel.
\item Scalable reactionless generators for clean, unlimited energy production.
\item Possible resolution of global energy constraints if laboratory tests scheduled for 2026--2029 confirm the predicted performance.
\end{itemize}

Energy conservation is preserved as the mechanism draws from the topological vacuum reservoir rather than creating energy from nothing.

\section{Conclusion}

Riemann sphere compactification of the supersymmetric Klein-bottle universe resolves infinite primordial knot cycles into a closed, unitary structure. The 3-6-9 invariant is derived rigorously from Euler's identity and the non-orientable topology. The framework provides concrete, falsifiable predictions and a consistent basis for local vacuum energy extraction.

\begin{thebibliography}{9}

\bibitem{soliman2025v1}
Soliman, S. (2025). 
The Supersymmetric Klein-Bottle Universe: 3-6-9 Invariant and Vacuum Entanglement (v1.0). 
\href{https://doi.org/10.5281/zenodo.17844202}{DOI: 10.5281/zenodo.17844202}

\bibitem{gibbons1977}
Gibbons, G. W., \& Hawking, S. W. (1977). 
Cosmological event horizons, thermodynamics, and particle creation. 
Phys. Rev. D 15, 2738.

\bibitem{verlinde2011}
Verlinde, E. (2011). 
On the origin of gravity and the laws of Newton. 
JHEP 04, 029.

\end{thebibliography}

\end{document}